\documentclass{article}

\usepackage{agda}
\usepackage{unicode}
\usepackage[utf8]{inputenc}

\title{Containers in Higher Kinds}

% Let lines with long inline code stretch instead of overflowing
\emergencystretch=1.5em

\begin{document}

\maketitle

Strictly positive functors can be represented as \emph{containers}~\cite{containers}: a type of shapes \AgdaBound{S}\;\AgdaSymbol{:}\;\AgdaPrimitive{Type} together with a family of positions \AgdaBound{P}\;\AgdaSymbol{:}\;\AgdaBound{S}\;\AgdaSymbol{→}\;\AgdaPrimitive{Type}, which gives rise to a functor \AgdaFunction{⟦}\;\AgdaBound{S}\;\AgdaOperator{\AgdaInductiveConstructor{◃}}\;\AgdaBound{P}\;\AgdaFunction{⟧}\;\AgdaSymbol{:}\;\AgdaPrimitive{Type}\;\AgdaSymbol{→}\;\AgdaPrimitive{Type}. Every container has an initial algebra, the W-type, and a terminal coalgebra, the M-type~\cite{containers,ahrens2015nonwellfounded}. However, some inductive and coinductive types do not fit the scheme of W-types and M-types. A typical example is the nested datatype~\cite{bird1998nested}:

\begin{code}%
\>[0]\AgdaKeyword{data}\AgdaSpace{}%
\AgdaDatatype{Bush}\AgdaSpace{}%
\AgdaSymbol{(}\AgdaBound{A}\AgdaSpace{}%
\AgdaSymbol{:}\AgdaSpace{}%
\AgdaPrimitive{Type}\AgdaSymbol{)}\AgdaSpace{}%
\AgdaSymbol{:}\AgdaSpace{}%
\AgdaPrimitive{Type}\AgdaSpace{}%
\AgdaKeyword{where}\<%
\\
\>[0][@{}l@{\AgdaIndent{0}}]%
\>[2]\AgdaInductiveConstructor{[]}\AgdaSpace{}%
\AgdaSymbol{:}\AgdaSpace{}%
\AgdaDatatype{Bush}\AgdaSpace{}%
\AgdaBound{A}\<%
\\
%
\>[2]\AgdaOperator{\AgdaInductiveConstructor{\AgdaUnderscore{}::\AgdaUnderscore{}}}\AgdaSpace{}%
\AgdaSymbol{:}\AgdaSpace{}%
\AgdaBound{A}\AgdaSpace{}%
\AgdaSymbol{→}\AgdaSpace{}%
\AgdaDatatype{Bush}\AgdaSpace{}%
\AgdaSymbol{(}\AgdaDatatype{Bush}\AgdaSpace{}%
\AgdaBound{A}\AgdaSymbol{)}\AgdaSpace{}%
\AgdaSymbol{→}\AgdaSpace{}%
\AgdaDatatype{Bush}\AgdaSpace{}%
\AgdaBound{A}\<%
\end{code}

Because of its nested recursion, \AgdaDatatype{Bush}\;\AgdaBound{A} cannot be obtained as a W-type. Nevertheless, \AgdaDatatype{Bush} itself is represented by a container \AgdaFunction{Bushᶜ}\;\AgdaSymbol{:}\;\AgdaRecord{Cont}, which arises as the initial algebra of a functor on the category of containers~\cite{bird1998nested,bird1999generalised,johann2007initial}. This motivates a notion of containers at higher kinds, such as \AgdaSymbol{(}\AgdaInductiveConstructor{*}\;\AgdaOperator{\AgdaInductiveConstructor{⇒}}\;\AgdaInductiveConstructor{*}\AgdaSymbol{)}\;\AgdaOperator{\AgdaInductiveConstructor{⇒}}\;\AgdaInductiveConstructor{*}\;\AgdaOperator{\AgdaInductiveConstructor{⇒}}\;\AgdaInductiveConstructor{*} or \AgdaSymbol{(}\AgdaInductiveConstructor{*}\;\AgdaOperator{\AgdaInductiveConstructor{⇒}}\;\AgdaInductiveConstructor{*}\AgdaSymbol{)}\;\AgdaOperator{\AgdaInductiveConstructor{⇒}}\;\AgdaInductiveConstructor{*}, capable of representing strictly positive functors on containers.

\paragraph{Second-Order Containers}

We first introduce a syntax \AgdaRecord{2Cont} of second-order containers. It extends ordinary containers with a family of positions for the function variable, \AgdaBound{PF}\;\AgdaSymbol{:}\;\AgdaBound{S}\;\AgdaSymbol{→}\;\AgdaPrimitive{Type}, and an inductive continuation \mbox{\AgdaBound{RF}\;\AgdaSymbol{:}\;\AgdaSymbol{(}\AgdaBound{s}\;\AgdaSymbol{:}\;\AgdaBound{S}\AgdaSymbol{)}\;\AgdaSymbol{→}\;\AgdaBound{PF}\;\AgdaBound{s}\;\AgdaSymbol{→}\;\AgdaRecord{2Cont}}, which describes the argument applied to the function variable. Every second-order container is interpreted as a second-order type-level function, \AgdaFunction{2⟦\_⟧}\;\AgdaSymbol{:}\;\AgdaRecord{2Cont}\;\AgdaSymbol{→}\;\AgdaSymbol{(}\AgdaPrimitive{Type}\;\AgdaSymbol{→}\;\AgdaPrimitive{Type}\AgdaSymbol{)}\;\AgdaSymbol{→}\;\AgdaPrimitive{Type}\;\AgdaSymbol{→}\;\AgdaPrimitive{Type}. In our example, the signature of \AgdaDatatype{Bush}: \AgdaFunction{SigBush}\;\AgdaBound{F}\;\AgdaBound{X}\;\AgdaSymbol{=}\;\AgdaDatatype{⊤}\;\AgdaOperator{\AgdaDatatype{⊎}}\;\AgdaBound{X}\;\AgdaOperator{\AgdaDatatype{×}}\;\AgdaBound{F}\;\AgdaSymbol{(}\AgdaBound{F}\;\AgdaBound{X}\AgdaSymbol{)}, is precisely the interpretation of a second-order container \AgdaFunction{SigBushᶜ}\;\AgdaSymbol{:}\;\AgdaRecord{2Cont}. Categorically, every second-order container \AgdaBound{H} gives rise to an endofunctor on the category of containers, \AgdaFunction{app}\;\AgdaBound{H}\;\AgdaSymbol{:}\;\AgdaRecord{Cont}\;\AgdaSymbol{→}\;\AgdaRecord{Cont}.

We also define a second-order least fixpoint operator \AgdaFunction{2W}\;\AgdaSymbol{:}\;\AgdaRecord{2Cont}\;\AgdaSymbol{→}\;\AgdaRecord{Cont} by induction-induction, and thereby recover \AgdaFunction{Bushᶜ} as the initial algebra of \AgdaFunction{app}\;\AgdaFunction{SigBushᶜ}. The second-order greatest fixpoint operator \AgdaFunction{2M}\;\AgdaSymbol{:}\;\AgdaRecord{2Cont}\;\AgdaSymbol{→}\;\AgdaRecord{Cont} is defined dually, using coinduction-induction.

\paragraph{Higher-Order Containers}

This story of constructing ``containers of containers'' can be iterated. We therefore introduce a general notion of higher-order containers, \AgdaFunction{HCont}\;\AgdaSymbol{:}\;\AgdaDatatype{Kind}\;\AgdaSymbol{→}\;\AgdaPrimitive{Type}, which covers containers of every kind. For brevity, we do not present the full syntax; we only remark that higher-order containers are closed normal forms of the simply-typed $\lambda$-calculus, defined mutually with neutral terms and spines, following the standard presentation of hereditary substitution~\cite{keller2010hereditary}. In this setting it is easy to see that \AgdaPrimitive{Type}\;\AgdaFunction{≅}\;\AgdaFunction{HCont}\;\AgdaInductiveConstructor{*}, \AgdaRecord{Cont}\;\AgdaFunction{≅}\;\AgdaFunction{HCont}\;\AgdaSymbol{(}\AgdaInductiveConstructor{*}\;\AgdaOperator{\AgdaInductiveConstructor{⇒}}\;\AgdaInductiveConstructor{*}\AgdaSymbol{)} and \AgdaRecord{2Cont}\;\AgdaFunction{≅}\;\AgdaFunction{HCont}\;\AgdaSymbol{((}\AgdaInductiveConstructor{*}\;\AgdaOperator{\AgdaInductiveConstructor{⇒}}\;\AgdaInductiveConstructor{*}\AgdaSymbol{)}\;\AgdaOperator{\AgdaInductiveConstructor{⇒}}\;\AgdaInductiveConstructor{*}\;\AgdaOperator{\AgdaInductiveConstructor{⇒}}\;\AgdaInductiveConstructor{*}\AgdaSymbol{)}.

Every higher-order container can be interpreted as a type-level function, \AgdaFunction{H⟦\_⟧}\;\AgdaSymbol{:}\;\AgdaFunction{HCont}\;\AgdaBound{A}\;\AgdaSymbol{→}\;\AgdaFunction{⟦}\;\AgdaBound{A}\;\AgdaFunction{⟧k}, where \AgdaFunction{⟦\_⟧k}\;\AgdaSymbol{:}\;\AgdaDatatype{Kind}\;\AgdaSymbol{→}\;\AgdaPrimitive{Type₁} is the standard interpretation of kinds: \AgdaFunction{⟦}\;\AgdaInductiveConstructor{*}\;\AgdaFunction{⟧k}\;\AgdaSymbol{=}\;\AgdaPrimitive{Type} and \AgdaFunction{⟦}\;\AgdaBound{A}\;\AgdaOperator{\AgdaInductiveConstructor{⇒}}\;\AgdaBound{B}\;\AgdaFunction{⟧k}\;\AgdaSymbol{=}\;\AgdaFunction{⟦}\;\AgdaBound{A}\;\AgdaFunction{⟧k}\;\AgdaSymbol{→}\;\AgdaFunction{⟦}\;\AgdaBound{B}\;\AgdaFunction{⟧k}.

Categorically, application of normal forms, computed by hereditary substitution,

\[
\AgdaFunction{napp}\;\AgdaSymbol{:}\;\AgdaFunction{HCont}\;\AgdaSymbol{(}\AgdaBound{A}\;\AgdaOperator{\AgdaInductiveConstructor{⇒}}\;\AgdaBound{B}\AgdaSymbol{)}\;\AgdaSymbol{→}\;\AgdaFunction{HCont}\;\AgdaBound{A}\;\AgdaSymbol{→}\;\AgdaFunction{HCont}\;\AgdaBound{B},
\]
turns every \AgdaBound{H}\;\AgdaSymbol{:}\;\AgdaFunction{HCont}\;\AgdaSymbol{(}\AgdaBound{A}\;\AgdaOperator{\AgdaInductiveConstructor{⇒}}\;\AgdaBound{B}\AgdaSymbol{)} into a functor \AgdaFunction{napp}\;\AgdaBound{H} from the category of higher-order containers of kind \AgdaBound{A} to that of kind \AgdaBound{B}, generalising \AgdaFunction{app} for second-order containers.

However, this semantics shows that there is no fixpoint operator beyond the second order. Consider the following third-order container, written as the corresponding $\lambda$-term since its encoding as a normal form is not very readable:

\begin{code}%
\>[0]\AgdaFunction{C}\AgdaSpace{}%
\AgdaSymbol{:}\AgdaSpace{}%
\AgdaFunction{HCont}\AgdaSpace{}%
\AgdaSymbol{(((}\AgdaInductiveConstructor{*}\AgdaSpace{}%
\AgdaOperator{\AgdaInductiveConstructor{⇒}}\AgdaSpace{}%
\AgdaInductiveConstructor{*}\AgdaSymbol{)}\AgdaSpace{}%
\AgdaOperator{\AgdaInductiveConstructor{⇒}}\AgdaSpace{}%
\AgdaInductiveConstructor{*}\AgdaSymbol{)}\AgdaSpace{}%
\AgdaOperator{\AgdaInductiveConstructor{⇒}}\AgdaSpace{}%
\AgdaSymbol{(}\AgdaInductiveConstructor{*}\AgdaSpace{}%
\AgdaOperator{\AgdaInductiveConstructor{⇒}}\AgdaSpace{}%
\AgdaInductiveConstructor{*}\AgdaSymbol{)}\AgdaSpace{}%
\AgdaOperator{\AgdaInductiveConstructor{⇒}}\AgdaSpace{}%
\AgdaInductiveConstructor{*}\AgdaSymbol{)}\<%
\\
\>[0]\AgdaFunction{C}\AgdaSpace{}%
\AgdaBound{F}\AgdaSpace{}%
\AgdaBound{G}\AgdaSpace{}%
\AgdaSymbol{=}\AgdaSpace{}%
\AgdaBound{G}\AgdaSpace{}%
\AgdaSymbol{(}\AgdaBound{F}\AgdaSpace{}%
\AgdaBound{G}\AgdaSymbol{)}\<%
\end{code}

\noindent
Suppose there were a third-order fixpoint operator
\[
  \AgdaFunction{3W}\;\AgdaSymbol{:}\;\AgdaFunction{HCont}\;\AgdaSymbol{(((}\AgdaInductiveConstructor{*}\;\AgdaOperator{\AgdaInductiveConstructor{⇒}}\;\AgdaInductiveConstructor{*}\AgdaSymbol{)}\;\AgdaOperator{\AgdaInductiveConstructor{⇒}}\;\AgdaInductiveConstructor{*}\AgdaSymbol{)}\;\AgdaOperator{\AgdaInductiveConstructor{⇒}}\;\AgdaSymbol{(}\AgdaInductiveConstructor{*}\;\AgdaOperator{\AgdaInductiveConstructor{⇒}}\;\AgdaInductiveConstructor{*}\AgdaSymbol{)}\;\AgdaOperator{\AgdaInductiveConstructor{⇒}}\;\AgdaInductiveConstructor{*}\AgdaSymbol{)}\;\AgdaSymbol{→}\;\AgdaFunction{HCont}\;\AgdaSymbol{((}\AgdaInductiveConstructor{*}\;\AgdaOperator{\AgdaInductiveConstructor{⇒}}\;\AgdaInductiveConstructor{*}\AgdaSymbol{)}\;\AgdaOperator{\AgdaInductiveConstructor{⇒}}\;\AgdaInductiveConstructor{*}\AgdaSymbol{)}.
\]
Then \AgdaFunction{Fix}\;\AgdaSymbol{=}\;\AgdaFunction{3W}\;\AgdaFunction{C} would satisfy \AgdaFunction{Fix}\;\AgdaBound{G}\;\AgdaFunction{≅}\;\AgdaBound{G}\;\AgdaSymbol{(}\AgdaFunction{Fix}\;\AgdaBound{G}\AgdaSymbol{)} for every \AgdaBound{G}; that is, \AgdaFunction{Fix} would be an \emph{internal} fixpoint operator. Under the semantics above, \AgdaFunction{H⟦}\;\AgdaFunction{Fix}\;\AgdaFunction{⟧}, a function of type \AgdaSymbol{(}\AgdaPrimitive{Type}\;\AgdaSymbol{→}\;\AgdaPrimitive{Type}\AgdaSymbol{)}\;\AgdaSymbol{→}\;\AgdaPrimitive{Type} would then yield a fixpoint of every type operator. Taking \AgdaBound{G}\;\AgdaBound{X}\;\AgdaSymbol{=}\;\AgdaBound{X}\;\AgdaSymbol{→}\;\AgdaDatatype{Bool} gives a type isomorphic to its own type of predicates, which is impossible by Cantor's diagonal argument.

We would like to incorporate initial algebras and terminal coalgebras at all kinds, either by adding explicit fixpoint operators to the syntax or by allowing infinite terms (i.e.\ using a coinductively defined syntax). In either case, the naive semantics sketched above has to be given up.

\paragraph{Mixing Higher-Order and Indexed}

It turns out that the higher-order structure is orthogonal to indexing, so our constructions combine freely with indexed containers~\cite{indexedcontainers}. We define a syntax of second-order indexed containers \AgdaKeyword{record}\;\AgdaRecord{2ICont}\;\AgdaSymbol{(}\AgdaBound{I}\;\AgdaBound{J}\;\AgdaSymbol{:}\;\AgdaPrimitive{Type}\AgdaSymbol{)}\;\AgdaSymbol{:}\;\AgdaPrimitive{Type₁}, with a type semantics
\[
  \AgdaFunction{2I⟦\_⟧}\;\AgdaSymbol{:}\;\AgdaRecord{2ICont}\;\AgdaBound{I}\;\AgdaBound{J}\;\AgdaSymbol{→}\;\AgdaSymbol{(}\AgdaBound{I}\;\AgdaSymbol{→}\;\AgdaPrimitive{Type}\;\AgdaSymbol{→}\;\AgdaPrimitive{Type}\AgdaSymbol{)}\;\AgdaSymbol{→}\;\AgdaSymbol{(}\AgdaBound{J}\;\AgdaSymbol{→}\;\AgdaPrimitive{Type}\;\AgdaSymbol{→}\;\AgdaPrimitive{Type}\AgdaSymbol{)},
\]
a functorial interpretation, and a least fixpoint given by an inductive family:
\[
  \begin{array}{l}
    \AgdaFunction{2Iapp}\;\AgdaSymbol{:}\;\AgdaRecord{2ICont}\;\AgdaBound{I}\;\AgdaBound{J}\;\AgdaSymbol{→}\;\AgdaSymbol{(}\AgdaBound{I}\;\AgdaSymbol{→}\;\AgdaRecord{Cont}\AgdaSymbol{)}\;\AgdaSymbol{→}\;\AgdaBound{J}\;\AgdaSymbol{→}\;\AgdaRecord{Cont}\\
    \AgdaKeyword{data}\;\AgdaDatatype{2IW}\;\AgdaSymbol{(}\AgdaBound{H}\;\AgdaSymbol{:}\;\AgdaRecord{2ICont}\;\AgdaBound{I}\;\AgdaBound{I}\AgdaSymbol{)}\;\AgdaSymbol{:}\;\AgdaBound{I}\;\AgdaSymbol{→}\;\AgdaPrimitive{Type}\;\AgdaSymbol{→}\;\AgdaPrimitive{Type}
  \end{array}
\]
For example, vectors \AgdaDatatype{Vec}\;\AgdaSymbol{(}\AgdaBound{X}\;\AgdaSymbol{:}\;\AgdaPrimitive{Type}\AgdaSymbol{)}\;\AgdaSymbol{:}\;\AgdaDatatype{ℕ}\;\AgdaSymbol{→}\;\AgdaPrimitive{Type}, indexed by their length, arise as an instance of \AgdaDatatype{2IW}. The same approach generalises to higher-order indexed containers.

\paragraph{Conclusion and Future Work}

We have formalised in Cubical Agda~\cite{cubicalagda} the syntax and semantics of second-order and higher-order containers, including the least and greatest fixpoints of second-order containers. This work is ongoing, and several questions remain open. In particular, we aim to show that W-types remain sufficient: for the second-order case, we conjecture that \AgdaFunction{2W} can be reduced to the first-order \AgdaDatatype{W}, analogously to the reduction of \AgdaFunction{2M} to \AgdaDatatype{M}. We also plan to develop a semantics of higher-order containers together with their fixpoints and recursors.

%% \begin{thebibliography}{9}

%% \bibitem{ahrens2015nonwellfounded}
%% Benedikt Ahrens, Paolo Capriotti, and R{\'e}gis Spadotti.
%% \newblock Non-wellfounded trees in homotopy type theory.
%% \newblock In \emph{13th International Conference on Typed Lambda Calculi and Applications (TLCA 2015)}, volume~38 of \emph{LIPIcs}, pages 17--30, 2015.

%% \bibitem{containers}
%% Michael Abbott, Thorsten Altenkirch, and Neil Ghani.
%% \newblock Containers: Constructing strictly positive types.
%% \newblock \emph{Theoretical Computer Science}, 342(1):3--27, 2005.

%% \bibitem{indexedcontainers}
%% Thorsten Altenkirch, Neil Ghani, Peter Hancock, Conor McBride, and Peter Morris.
%% \newblock Indexed containers.
%% \newblock \emph{Journal of Functional Programming}, 25:e5, 2015.

%% \bibitem{bird1998nested}
%% Richard Bird and Lambert Meertens.
%% \newblock Nested datatypes.
%% \newblock In \emph{Mathematics of Program Construction (MPC 1998)}, volume 1422 of \emph{LNCS}, pages 52--67. Springer, 1998.

%% \bibitem{bird1999generalised}
%% Richard Bird and Ross Paterson.
%% \newblock Generalised folds for nested datatypes.
%% \newblock \emph{Formal Aspects of Computing}, 11(2):200--222, 1999.

%% \bibitem{johann2007initial}
%% Patricia Johann and Neil Ghani.
%% \newblock Initial algebra semantics is enough!
%% \newblock In \emph{Typed Lambda Calculi and Applications (TLCA 2007)}, volume 4583 of \emph{LNCS}, pages 207--222. Springer, 2007.

%% \bibitem{keller2010hereditary}
%% Chantal Keller and Thorsten Altenkirch.
%% \newblock Hereditary substitutions for simple types, formalized.
%% \newblock In \emph{Proceedings of the 3rd Workshop on Mathematically Structured Functional Programming (MSFP 2010)}, pages 3--10. ACM, 2010.

%% \bibitem{cubicalagda}
%% Andrea Vezzosi, Anders M{\"o}rtberg, and Andreas Abel.
%% \newblock Cubical {Agda}: A dependently typed programming language with univalence and higher inductive types.
%% \newblock \emph{Proceedings of the ACM on Programming Languages}, 3(ICFP):87:1--87:29, 2019.

%% \end{thebibliography}

\end{document}
